
\section{Cauchy-Binet Formula}

Application 1: {\bf Example Partitioned Matrix and the Cauchy-Binet Formula}
\begin{equation}
M =
\left[
\begin{array}{cccc|cccc}
0 & 0 & \hdots & 0 & b_{11} & b_{12} & \hdots & b_{1n}\\
0 & 0 & \hdots & 0 & b_{21} & b_{22} & \hdots & b_{2n}\\
\hdotsfor[2]{3} & \hdots &\hdotsfor[2]{3} & \hdots\\
0 & 0 & \hdots & 0 & b_{m1} & b_{m2} & \hdots & b_{mn}\\
\hline
c_{11} & c_{12} &\hdots & c_{1m} & 1 & 0& \hdots & 0\\
c_{21} & c_{22} &\hdots & c_{2m} & 0 & 1& \hdots & 0\\
\hdotsfor[2]{3} & \hdots & \hdotsfor[2]{3} & 0\\
c_{n1} & c_{n2} &\hdots & c_{nm} & 0 & 0& \hdots& 1
\end{array} \right]
\end{equation}
In the above $M$ is an ($m+n$)-by-($m+n$) matrix which is partitioned into: $A$ an $m$-by-$m$ matrix of 0's, $B$ an $m$-by-$n$ matrix, $C$ an $n$-by-$m$ matrix and $D$ is the $n$-by-$n$ Identity matrix. Then using the determinant in the Schur's Complement form we have
$$\det M = \det D \det(A - BD^{-1}C) = \det (-BC) = (-1)^m \det(BC).$$\\

We can obtain a similar result by row operations. Denote the ($m+i)$ row by $\a_i$. Adding a multiple of row $i$ to row $j$ does not change the value of $\det M$.
We successively subtract the $\sum_{i=1}^n b_{k\, i} \a_i$ from the $k$-th row of M, for $k=1, 2, ..., m$ changing M into a different matrix with the same determinant, viz.,
\begin{equation}
\tilde M=
\left[
\begin{array}{cccc|cccc}
-d_{11} & -d_{12} & \hdots & -d_{1m} & 0 & 0& \hdots & 0\\
-d_{21} & -d_{22} & \hdots & -d_{2m} & 0 & 0 & \hdots & 0\\
\hdotsfor[2]{3} & \hdots &\hdotsfor[2]{3} & \hdots\\
-d_{m1} & -d_{m2} & \hdots & -d_{mm} & 0 & 0 & \hdots & 0\\
\hline
c_{11} & c_{12} &\hdots & c_{1m} & 1 & 0& \hdots & 0\\
c_{21} & c_{22} &\hdots & c_{2m} & 0 & 1& \hdots & 0\\
\hdotsfor[2]{3} & \hdots & \hdotsfor[2]{3} & 0\\
c_{n1} & c_{n2} &\hdots & c_{nm} & 0 & 0& \hdots& 1
\end{array} \right],
\end{equation}where $\displaystyle{d_{kj} = \sum_{i=1}^n \, b_{ki} c_{ij}}, ~\mbox{ for } k,j = 1\,,...,\,m$. Since $\det M = \det \tilde M$ we have another demonstration of
\begin{equation}
\det M = (-1)^m \det (BC)
\end{equation} where $B$ and $C$ are the matrices corresponding to $b_{kj}$ and $c_{ij}$ respectively.\\

Another representation of $\det M$ can be obtained by applying Laplace's Generalized Theorem successively. First swap the bottom $n$ rows with the top $m$ rows (by moving each of the bottom $n$ rows up $m$ positions and changing the sign of $\det M$ accordingly
\begin{equation}
\det M = (-1)^{m^2}
\det \left[
\begin{array}{cccc|cccc}
c_{11} & c_{12} &\hdots & c_{1m} & 1 & 0& \hdots & 0\\
c_{21} & c_{22} &\hdots & c_{2m} & 0 & 1& \hdots & 0\\
\hdotsfor[2]{3} & \hdots & \hdotsfor[2]{3} & 0\\
c_{n1} & c_{n2} &\hdots & c_{nm} & 0 & 0& \hdots& 1 \\
\hline
0 & 0 & \hdots & 0 & b_{11} & b_{12} & \hdots & b_{1n}\\
0 & 0 & \hdots & 0 & b_{21} & b_{22} & \hdots & b_{2n}\\
\hdotsfor[2]{3} & \hdots &\hdotsfor[2]{3} & \hdots\\
0 & 0 & \hdots & 0 & b_{m1} & b_{m2} & \hdots & b_{mn}
\end{array} \right]
\end{equation}

Using Laplace's Generalized Theorem the first time across the bottom $m$ rows and taking advantage of the $m$-by-$m$ block of zeros in the lower left hand corner to drop some of the minors in the expansion we have only a sum over the $m$-by-$m$ minors in the $b_{ij}$ elements starting in columns $m+1$ through $m+n$ of $M$. After forming any one of these minors we have to select $m$ columns to remove to form the algebraic complement to use in the Laplace Theorem. The $n$-by-$n$ complementary minor consists of the $n$-by-$m$ $c_{ij}$ matrix elements with an additional $n-m$ columns which remain in from the former Identity matrix in the upper right hand side of the working matrix.\\

We now apply Laplace's Generalized Theorem a second time, now expanding over the first $m$ columns of the complementary $n$-by$n$ minor. We must again sum over the choice of $m$ rows from the first $n$ rows. The complementary minor to this second set of minors is obtained from the original $n$-by-$n$ Identity matrix, $I_n$, with $m$ columns removed by forming the first set of minors in the $b_{ij} $ and removing the $m$ rows to form the second set of minors. We are left with an $n-m$-by-$n-m$ matrix from $I_{n}$. If the rows that were struck out do not match exactly the columns that were struck out then there will appear in the the reduced matrix rows and columns with all zero entries and the corresponding determinant vanishes. So in the second application of Laplace's Generalized theorem there is only one term surviving in the sum, for any given choice of $m$ numbers
$1 \le \rho_1 \le \rho_2 \le ...\le \rho_m \le n$, the only term that survive is
\begin{equation}
\begin{vmatrix}
c_{1\rho_1} & c_{1\rho_2} & \hdots & c_{1\rho_m} \\
c_{2\rho_1} & c_{2\rho_2} & \hdots & c_{2\rho_m} \\
\hdotsfor[2]{4}\\
c_{m\rho_1} & c_{m\rho_2} & \hdots & c_{m\rho_m} \\
\end{vmatrix}
\begin{vmatrix}
b_{1\rho_1} & b_{1\rho_2} & \hdots & b_{1\rho_m} \\
b_{2\rho_1} & b_{2\rho_2} & \hdots & b_{2\rho_m} \\
\hdotsfor[2]{4}\\
b_{m\rho_1} & b_{m\rho_2} & \hdots & b_{m\rho_m} \\
\end{vmatrix}
\end{equation}

Putting both results together and using the fact that both methods give an expression for $\det M$ we obtain the Cauchy-Binet Formula for matrices $B$ ($n$-by-$m$) and $C$ ($m$-by-$n$) with $m\le n$:
\begin{equation}
\det(BC) =\sum_{\substack{\rho_1, \rho_2, ...,\rho_m =1 \\ \rho_1 < \rho_2 < ... < \rho_m}}^n \,
\begin{vmatrix}
c_{1\rho_1} & c_{1\rho_2} & \hdots & c_{1\rho_m} \\
c_{2\rho_1} & c_{2\rho_2} & \hdots & c_{2\rho_m} \\
\hdotsfor[2]{4}\\
c_{m\rho_1} & c_{m\rho_2} & \hdots & c_{m\rho_m} \\
\end{vmatrix}
\begin{vmatrix}
b_{1\rho_1} & b_{1\rho_2} & \hdots & b_{1\rho_m} \\
b_{2\rho_1} & b_{2\rho_2} & \hdots & b_{2\rho_m} \\
\hdotsfor[2]{4}\\
b_{m\rho_1} & b_{m\rho_2} & \hdots & b_{m\rho_m} \\
\end{vmatrix}
\end{equation}\\

{\bf 1: For rectangular matrices, when the number of columns exceeds the number of rows.} The overall determinant is formed summing the product of all the sub-determinants be formed from one array {(\it by taking a number of columns equal to the number of rows}) multiplied by the corresponding determinants from the other array.\\

{\bf 2: When the number of rows exceeds the number of columns, the resulting determinant vanishes.} It will be observed that the resulting determinant is the same as would arise if an appropriate number of zero-columns were added to each of the given arrays and then forming the determinant of the product.\\

Application 2: {\bf Alternate Derivation of the Cauchy-Binet Formula}\\
Let $A$ be an $m$-by-$n$ matrix and let $B$ be an $n$-by-$m$ matrix. Let $1\le j_1,j_2,...,j_m\le n$.
Let $A_{j_1j_2...j_m}$ denote the $m$-by-$m$ matrix consisting of the $m$ columns $j_1,j_2,...,j_m$ of $A$.
Let $B_{j_1j_2...j_m}$ denote the $m$-by-$m$ matrix consisting of the $m$ rows $j_1,j_2,...,j_m$ of $B$.
Let $\{k_1,k_2,...,k_m\}$ be an ordered $m$-tuple of integers and let $\{j_1,j_2,...,j_m\}$ be the same set of integers but arranged into non-decreasing order:
$$\j_1\le j_2 \le ... \le j_m.$$
Then we have the relationships that
\begin{eqnarray*}
\det \left(B_{k_1k_2...k_m}\right) &=& \epsilon_{k_1k_2...k_m} \det \left(B_{j_1j_2...j_m}\right),\\
\det \left(B_{j_1j_2...j_m} \right) &=& \epsilon_{k_1k_2...k_m} \det \left(B_{k_1k_2...k_m}\right)\mbox{ no implied sum on } k_i's.
\end{eqnarray*}\\
Now from the definition of determinant for the $m$-by-$m$ matrix $AB$ we have
\begin{eqnarray*}
\det (AB) &=& \sum_{1\le l_1,...l_m \le m} \, \epsilon_{l_1l_2...l_m} \left( \sum_{k=1}^n a_{1k} b_{kl_1} \right) ... \left( \sum_{k=1}^n a_{mk}b_{kl_m}\right) \\
&=& \sum_{1\le k_1,...k_m \le n} \, \left(a_{1k_1}...a_{mk_m}\right) \sum_{1\le l_1,...l_m \le m} \, \epsilon_{l_1l_2...l_m} b_{k_1l_1} ... b_{k_ml_m} \\
&=& \sum_{1\le k_1,...k_m \le n} \, \left( a_{1k_1}...a_{mk_m} \right) \det \left(B_{k_1...k_m}\right)\\
&=& \sum_{1\le k_1,...k_m \le n} \, \left(a_{1k_1}...a_{mk_m} \epsilon_{k_1...k_m}\right) \det \left(B_{j_1...j_m}\right)\\
&=& \sum_{1\le k_1,...k_m \le n} \, \left(\epsilon_{k_1...k_m} \, a_{1k_1}...a_{mk_m}\right) \det \left(B_{j_1...j_m}\right)\\
&=& \sum_{1\le j_1 \le j_2 \le ...\le j_m \le n} \det \left(A_{j_1...j_m}\right)\,\det \left(B_{j_1...j_m}\right) \\
\end{eqnarray*}
If any two $j$'s are equal then $\det \left(A_{j_1...j_m}\right) = 0$\\
In the case of $A$ and $B = A^T$ we immediately have
$$ \det (AA^T) = \sum_{1\le j_1 \le j_2 \le ...\le j_m \le n} \left[ \det \left(A_{j_1...j_m}\right)\right]^2 $$
